\documentclass[10pt]{article}
\usepackage[T1]{fontenc}
\usepackage[utf8]{inputenc}
\usepackage[french]{babel}
\usepackage{amsmath,amssymb,mathrsfs,mathpazo}
\usepackage{geometry}
\geometry{a4paper,margin=2cm}
\usepackage{graphicx,array,tabularx,tabularray,colortbl,multirow,multicol,booktabs,enumitem,pifont,ulem,fancybox,diagbox,cancel}
\usepackage{pstricks,pstricks-add,pst-plot,pst-tree,pst-node,pst-text,pst-func,pst-math,pst-bspline,pst-3dplot,pst-eucl,pst-solides3d}
\let\pscalculate\undefined
\usepackage{pst-bezier}
\usepackage{tikz,pgfplots,tkz-tab}
\usetikzlibrary{arrows,arrows.meta,patterns,calc,positioning,trees,decorations.markings,shapes.geometric,shapes.misc,shapes.arrows,intersections,decorations.pathmorphing,decorations.pathreplacing,backgrounds,shadows,babel}
\pgfplotsset{compat=1.18}
\usepgfplotslibrary{fillbetween}
\usepackage[np]{numprint}
\usepackage{listings,makecell,pas-tableur,twoopt}
\let\barre\undefined
\usepackage[pstricks]{tablvar}
\definecolor{darkgreen}{RGB}{0,100,0}
\providecommand{\red}{\color{red}}
\providecommand{\blue}{\color{blue}}
\providecommand{\green}{\color{green}}
\providecommand{\black}{\color{black}}
\providecommand{\white}{\color{white}}
\providecommand{\gray}{\color{gray}}
\providecommand{\cyan}{\color{cyan}}
\providecommand{\magenta}{\color{magenta}}
\providecommand{\yellow}{\color{yellow}}
\providecommand{\violet}{\color{violet}}
\providecommand{\orange}{\color{orange}}
\providecommand{\euro}{\texteuro}
\providecommand{\EUR}{\texteuro}
\providecommand{\ieme}{\textsuperscript{e}}
\providecommand{\degres}{\ensuremath{^\circ}}
\providecommand{\degree}{\ensuremath{^\circ}}
\providecommand{\vv}[1]{\overrightarrow{#1}}
\providecommand{\Coord}[1]{#1}
\providecommand{\decosix}{}
\providecommand{\starredbullet}{\ensuremath{\bullet}}
\pagestyle{empty}
\setlength{\parindent}{0pt}
\usepackage{adjustbox}\DeclareUnicodeCharacter{2236}{:}\DeclareUnicodeCharacter{2190}{\ensuremath{\leftarrow}}
\providecommand{\R}{\mathbb{R}}
\providecommand{\N}{\mathbb{N}}
\providecommand{\Z}{\mathbb{Z}}
\providecommand{\Q}{\mathbb{Q}}
\providecommand{\D}{\mathbb{D}}
\providecommand{\C}{\mathbb{C}}
\providecommand{\e}{{\mathrm{e}}}
\providecommand{\vect}[1]{\overrightarrow{#1}}
\providecommand{\vectt}[1]{\overrightarrow{\mathrm{#1}}}
\providecommand{\vvt}[1]{\overrightarrow{\mathrm{#1}}}
\providecommand{\np}[2][]{{#2~#1}}
\providecommand{\vv}[1]{{\overrightarrow{#1}}}
\renewcommand{\d}{{\mathrm{d}}}
\providecommand{\degres}{{^\circ}}
\providecommand{\degre}{{^\circ}}
\providecommand{\blue}{\color{blue}}
\providecommand{\red}{\color{red}}
\providecommand{\enskip}{\quad}
\providecommand{\llbracket}{\lbrack\!\lbrack}
\providecommand{\rrbracket}{\rbrack\!\rbrack}
\providecommand{\euro}{{\text{€}}}
\providecommand{\no}{{\text{n°}}}
\providecommand{\interval}[2]{{[#1;#2]}}
\providecommand{\intervFO}[2]{{[#1;#2[}}
\providecommand{\dotfill}{\ldots}

\providecommand{\ds}{\displaystyle}
\providecommand{\textding}[1]{#1}
\providecommand{\Oij}{\(\mathrm{O};\vec{\imath},\vec{\jmath}\)}
\providecommand{\Oijk}{\(\mathrm{O};\vec{\imath},\vec{\jmath},\vec{k}\)}
\providecommand{\Ouv}{\(\mathrm{O};\vec{u},\vec{v}\)}
\providecommand{\intervFF}[2]{\ensuremath{[#1;#2]}}
\renewcommand{\euro}{\ifmmode\text{\texteuro}\else\texteuro\fi}
\DeclareUnicodeCharacter{00B0}{\ifmmode^{\circ}\else\textdegree\fi}
\DeclareUnicodeCharacter{20AC}{\ifmmode\text{\texteuro}\else\texteuro\fi}
% Énoncés : Éducation nationale. Transcription et corrigés : APMEP, Denis Vergès et contributeurs. Découpage et classement : CoopMaths, Sébastien Lozano, Éric Elter et contributeurs. Adaptation de la présentation : Maths Explorer.
% CC BY-SA : https://coopmaths.fr/www/about/
% https://coopmaths.fr/alea/static/bac/2022/tex/bac_2022_10_sujet1_caledonie_2.tex
\begin{document}
\noindent\textit{APMEP / CoopMaths. CC BY-SA. Présentation adaptée par Maths Explorer.}\par\medskip
\medskip

\textbf{Principaux domaines abordés :}
suites ;
fonctions, fonction exponentielle.

\medskip

On considère la fonction $f$ définie sur $\R$ par

\[f(x) = x^3\e^x.\]

On admet que la fonction $f$ est dérivable sur $\R$ et on note $f'$ sa fonction dérivée.

\medskip

\begin{enumerate}
\item On définit la suite $\left(u_n\right)$ par $u_0 = - 1$ et, pour tout entier naturel $n,\: u_{n+1} = f\left(u_n\right)$.
	\begin{enumerate}
		\item Calculer $u_1$ puis $u_2$.
		
On donnera les valeurs exactes, puis les valeurs approchées à $10^{-3}$.
		\item On considère la fonction \texttt{fonc}, écrite en langage Python ci-dessous.

\smallskip

\begin{minipage}{0.524\linewidth}
On rappelle qu'en langage Python, 

\og \texttt{i in range (n)}\fg{} signifie que 

\texttt{i} varie de 0 à \texttt{n -1}.
\end{minipage}\hfill
\begin{minipage}{0.40\linewidth}
\begin{tabular}{|l l|}\hline
\texttt{def}& \texttt{fonc (n)} :\\
&\texttt{u =- 1}\\
&\texttt{for i in range(n) :}\\
&\quad \texttt{u=u**3*exp(u)}\\
&\texttt{return u}\\\hline
\end{tabular}
\end{minipage}

Déterminer, sans justifier, la valeur renvoyée par \texttt{fonc (2)} arrondie à $10^{-3}$.
	\end{enumerate}
\item 
	\begin{enumerate}
		\item Démontrer que, pour tout $x$ réel, on a $f'(x) = x^2\e^x(x + 3)$.
		\item Justifier que le tableau de variations de $f$ sur $\R$ est celui représenté
ci-dessous :

\begin{center}
\psset{unit=1cm,arrowsize=2pt 3}
\begin{pspicture}(7,2.5)
\psframe(7,2.5)
\psline(0,2)(7,2)\psline(1,0)(1,2.5)
\uput[u](0.5,1.9){$x$} \uput[u](1.2,1.9){$- \infty$} \uput[u](4,1.9){$-3$} \uput[u](6.5,1.9){$+ \infty$} 
\rput(0.5,1){$f$}\uput[d](1.2,2){0}\uput[u](4,0){$- 27\e^{-3}$}\uput[d](6.5,2){$+ \infty$}
\psline{->}(1.5,1.5)(3.5,0.5)\psline{->}(4.5,0.52)(6.5,1.5)
\end{pspicture}
\end{center}

		\item Démontrer, par récurrence, que pour tout entier naturel $n$, on a :
		
\[- 1 \leqslant u_n \leqslant u_{n+1} \leqslant 0.\]
		
		\item En déduire que la suite $\left(u_n\right)$ est convergente.
		\item On note $\ell$ la limite de la suite $\left(u_n\right)$.

On rappelle que $\ell$ est solution de l'équation $f(x) = x$.

Déterminer $\ell$. (Pour cela, on admettra que l'équation $x^2\e^x - 1 = 0$ possède une seule solution dans $\R$ et que celle-ci est strictement supérieure à $\dfrac12$).
	\end{enumerate}
\end{enumerate}

\bigskip

\end{document}
